\documentclass[11pt, a4paper]{article}

% bfw-style.tex — gemeinsame Style-Definitionen für alle Handouts/Aufgabenhefte
% Voraussetzung: Kompilation mit xelatex (wegen fontspec)

\usepackage[ngerman]{babel}
\usepackage{geometry}
\geometry{a4paper, top=2.2cm, bottom=2.2cm, left=2.3cm, right=2.3cm}

\usepackage{fontspec}
% Fallback-Kette: Source Sans Pro -> Calibri -> Segoe UI -> Latin Modern Sans
% (Latin Modern ist immer mit MiKTeX vorhanden)
\IfFontExistsTF{Source Sans Pro}{%
  \setmainfont{Source Sans Pro}[Ligatures=TeX]%
}{\IfFontExistsTF{Calibri}{%
  \setmainfont{Calibri}[Ligatures=TeX]%
}{\IfFontExistsTF{Segoe UI}{%
  \setmainfont{Segoe UI}[Ligatures=TeX]%
}{%
  \setmainfont{Latin Modern Sans}[Ligatures=TeX]%
}}}
\IfFontExistsTF{Source Code Pro}{%
  \setmonofont{Source Code Pro}[Scale=0.92]%
}{\IfFontExistsTF{Consolas}{%
  \setmonofont{Consolas}[Scale=0.92]%
}{%
  \setmonofont{Latin Modern Mono}[Scale=0.92]%
}}

\usepackage{xcolor}
% Farbpalette: hell, freundlich, modern
\definecolor{bfwPrimary}{HTML}{0EA5A4}    % Petrol/Türkis - Primär
\definecolor{bfwPrimaryDark}{HTML}{0F766E} % dunkler Petrol für Headings
\definecolor{bfwAccent}{HTML}{F59E0B}     % Amber - Akzent
\definecolor{bfwSuccess}{HTML}{10B981}    % Grün - Erfolg / Lösung
\definecolor{bfwWarn}{HTML}{F97316}       % Orange - Achtung
\definecolor{bfwInk}{HTML}{0F172A}        % fast Schwarz
\definecolor{bfwInkSoft}{HTML}{475569}    % gedämpftes Grau
\definecolor{bfwBgSoft}{HTML}{F8FAFC}     % off-white
\definecolor{bfwBgTeal}{HTML}{ECFEFF}     % helles Türkis
\definecolor{bfwBgAmber}{HTML}{FEF3C7}    % helles Gelb
\definecolor{bfwBgRose}{HTML}{FFE4E6}     % helles Rosé für Eselsbrücken

\usepackage[colorlinks=true, linkcolor=bfwPrimaryDark, urlcolor=bfwPrimaryDark]{hyperref}
\usepackage{enumitem}
\setlist[itemize]{leftmargin=*, itemsep=2pt, topsep=2pt}
\setlist[enumerate]{leftmargin=*, itemsep=2pt, topsep=2pt}

\usepackage[most]{tcolorbox}
\usepackage{booktabs}
\usepackage{tikz}
\usetikzlibrary{decorations.pathmorphing, positioning, shapes.geometric, arrows.meta}

% Merkkästchen — wichtige Definition / Kernpunkt
\newtcolorbox{merksatz}[1][]{%
  enhanced, breakable,
  colback=bfwBgTeal,
  colframe=bfwPrimary,
  boxrule=0.6pt,
  arc=4pt,
  left=10pt, right=10pt, top=8pt, bottom=8pt,
  fonttitle=\bfseries\color{bfwPrimaryDark},
  title={\faLightbulb\ Merksatz},
  #1
}

% Eselsbrücke — Mnemoniks
\newtcolorbox{eselsbruecke}[1][]{%
  enhanced, breakable,
  colback=bfwBgRose,
  colframe=bfwAccent,
  boxrule=0.6pt,
  arc=4pt,
  left=10pt, right=10pt, top=8pt, bottom=8pt,
  fonttitle=\bfseries\color{bfwWarn},
  title={Eselsbrücke},
  #1
}

% Beispiel-Box
\newtcolorbox{beispiel}[1][]{%
  enhanced, breakable,
  colback=bfwBgSoft,
  colframe=bfwInkSoft,
  boxrule=0.4pt,
  arc=3pt,
  left=10pt, right=10pt, top=8pt, bottom=8pt,
  fonttitle=\bfseries\color{bfwInk},
  title={Beispiel},
  #1
}

% Lösungs-Box (für Aufgabenhefte)
\newtcolorbox{loesung}[1][]{%
  enhanced, breakable,
  colback=bfwBgAmber!40,
  colframe=bfwSuccess,
  boxrule=0.6pt,
  arc=3pt,
  left=10pt, right=10pt, top=8pt, bottom=8pt,
  fonttitle=\bfseries\color{bfwSuccess!70!black},
  title={Lösung},
  #1
}

% Glossar-Eintrag
\newcommand{\glossareintrag}[2]{%
  \par\noindent\textbf{\color{bfwPrimaryDark}#1} \enspace #2\par\smallskip
}

% Section-Stil — etwas farbiger als Standard
\usepackage{titlesec}
\titleformat{\section}
  {\Large\bfseries\color{bfwPrimaryDark}}
  {\thesection}{0.6em}{}
\titleformat{\subsection}
  {\large\bfseries\color{bfwPrimary}}
  {\thesubsection}{0.5em}{}
\titleformat{\subsubsection}
  {\normalsize\bfseries\color{bfwInk}}
  {\thesubsubsection}{0.4em}{}

% TikZ-Stil "skizzenhaft" — etwas weicher als Rough.js, aber stabil
\tikzset{
  roughbox/.style = {
    draw=bfwPrimaryDark, thick, fill=bfwBgTeal,
    rounded corners=4pt, inner sep=6pt
  },
  roughaccent/.style = {
    draw=bfwAccent, thick, fill=bfwBgAmber,
    rounded corners=4pt, inner sep=6pt
  },
  roughline/.style = {
    draw=bfwInkSoft, thick, -{Stealth[length=2.5mm]}
  }
}

% Bullet-Symbole (bewusst kein FontAwesome — vermeidet Font-Installations-Probleme)
\newcommand{\faLightbulb}{$\bullet$}
\newcommand{\faBookmark}{$\bullet$}

% Header/Footer
\usepackage{fancyhdr}
\pagestyle{fancy}
\fancyhf{}
\renewcommand{\headrulewidth}{0pt}
\fancyfoot[L]{\small\color{bfwInkSoft}\thepage}
\fancyfoot[R]{\small\color{bfwInkSoft} BFW M-064 Beschaffung im Büromanagement}


\title{\vspace*{-1.5cm}\Huge\bfseries\color{bfwPrimaryDark}Tag~4: Ausführen von Bestellungen\\[0.4em]
  \large\color{bfwInkSoft}\textnormal{Bestellverfahren, Andler-Formel und Bestellüberwachung}}
\author{\textsf{Beschaffung im Büromanagement}\\
\textsf{\small\copyright\ Can Siebert\,\textperiodcentered\,2026}}
\date{15.05.2026}

\begin{document}
\maketitle
\thispagestyle{empty}

\vspace{1em}
\begin{merksatz}[title={\bfseries Worum geht es heute?}]
Heute klären wir die zwei zentralen Fragen des operativen Einkaufs: \emph{Wann
bestellen?} und \emph{Wie viel bestellen?} Dazu betrachten wir vier Bestellverfahren,
die berühmte Andler-Formel zur optimalen Bestellmenge und das systematische Vorgehen
bei Lieferverzug. Die Andler-Formel ist eine der häufigsten Klausuraufgaben in der
IHK-Prüfung Teil~2 --- diese sollten Sie sicher rechnen können.
\end{merksatz}

\tableofcontents
\vspace{1em}
\hrule
\vspace{1em}

\section{Bestellverfahren}

\subsection{Die zwei Hauptgruppen}

In der Praxis gibt es vier Bestellverfahren, die sich in zwei Hauptgruppen einteilen
lassen:

\textbf{Verbrauchsgesteuerte Verfahren} reagieren auf den tatsächlichen Lagerverbrauch.
Hier zählen das Bestellpunktverfahren und JIT/Kanban.

\textbf{Zeitgesteuerte Verfahren} bestellen in festen Intervallen. Hierzu gehört das
Bestellrhythmusverfahren.

\subsection{Bestellpunktverfahren}

Bei diesem Verfahren wird eine Bestellung ausgelöst, sobald der Lagerbestand einen
festgelegten Wert --- den \emph{Meldebestand} --- erreicht.

\begin{merksatz}[title={Formel}]
\textbf{Meldebestand = Mindestbestand + (Tagesverbrauch $\times$ Lieferzeit in Tagen)}
\end{merksatz}

Drei Bestände muss man unterscheiden:

\begin{itemize}[itemsep=0pt]
  \item \textbf{Mindestbestand}: Sicherheitspuffer; sollte nie unterschritten werden.
  \item \textbf{Meldebestand}: Bestellauslösepunkt --- Mindestbestand plus erwarteter Verbrauch während der Lieferzeit.
  \item \textbf{Höchstbestand}: Mindestbestand + Bestellmenge.
\end{itemize}

\begin{beispiel}
Eine Druckerei verbraucht pro Werktag durchschnittlich 50~Blatt Premium-Papier.
Die Lieferzeit beträgt 5~Werktage. Der Sicherheits\-bestand liegt bei 100~Blatt.

\medskip
$\text{Meldebestand} = 100 + (50 \times 5) = \textbf{350 Blatt}$

\medskip
Sobald der Lagerbestand auf 350 Blatt fällt, wird neu bestellt. Während der 5-tägigen
Lieferzeit verbraucht die Druckerei weitere $50 \times 5 = 250$~Blatt $\to$ Bestand
sinkt auf 100~Blatt (genau auf den Sicherheits\-bestand). Die neue Lieferung kommt
just zum richtigen Zeitpunkt an.
\end{beispiel}

\subsection{Bestellrhythmusverfahren}

Beim Bestellrhythmus­verfahren wird in festen Zeit­intervallen bestellt --- z.\,B.\ jeden
Montag oder am 1.\ jedes Monats. Die Bestellmenge ist variabel und richtet sich danach,
wie viel der Bestand bis zum gewünschten Höchst\-bestand wieder aufgefüllt werden muss.

\begin{center}
\begin{tabular}{p{4cm} p{4.5cm} p{4.5cm}}
\toprule
\textbf{Kriterium} & \textbf{Bestellpunkt} & \textbf{Bestellrhythmus}\\
\midrule
Auslöser & Bestand erreicht Meldebestand & Festgelegter Zeitpunkt\\
Bestellmenge & Konstant & Variabel (auffüllen bis Höchstbestand)\\
Überwachung & Permanent & Nur zum Bestelltermin\\
Geeignet für & A-Teile, hoher Wert & B-/C-Teile, geringerer Wert\\
\bottomrule
\end{tabular}
\end{center}

\subsection{Just-in-Time und Kanban}

\textbf{Just-in-Time (JIT)} ist die fertigungssynchrone Beschaffung: Material wird
genau zum Bedarfszeit\-punkt geliefert. Lager praktisch null, Kapital­bindung minimal.
Voraussetzungen: sehr zuverlässige Lieferanten, kurze und planbare Lieferzeiten,
enge IT-Integration (oft EDI), oft auch geographische Nähe. Hauptrisiko: Lieferketten\-
störungen treffen die Produktion direkt --- wie die Halbleiterknappheit 2021/22 zeigte,
als Automobil­werke wegen ausbleibender Mikrochips tagelang stillstanden.

\textbf{Kanban} ist eine Pull-Variante des JIT, die in den 1970er Jahren bei Toyota
entwickelt wurde. Material wird erst geliefert, wenn der Verbraucher es signalisiert
--- traditionell über eine Karte (\emph{kanban} = jap.\ Karte), heute oft elektronisch.
Vorteil: Das System steuert sich selbst, ohne zentrale Planung.

\begin{eselsbruecke}
\textbf{,,Push vs.\ Pull''}: Klassische Beschaffung pusht Material in das Lager
(„auf Vorrat''), JIT/Kanban zieht es bedarfsgerecht aus dem Lieferanten heraus.
\end{eselsbruecke}

\section{Optimale Bestellmenge --- die Andler-Formel}

\subsection{Der Zielkonflikt}

Wer viel auf einmal bestellt, spart Bestellkosten und hat oft Mengenrabatte ---
zahlt aber hohe Lagerkosten. Wer wenig auf einmal bestellt, hat niedrige Lagerkosten,
aber viele teure Bestellvorgänge. Die optimale Bestellmenge ist die Menge, bei der die
\emph{Summe} aus Lager- und Bestell­kosten am niedrigsten ist.

\begin{itemize}[itemsep=0pt]
  \item \emph{Bestellkosten} fallen pro Bestellvorgang an: Personal für Anfrage und Bestellung, IT, Frachtkosten-Sockel. Sie sinken pro Stück, je größer die Bestellmenge.
  \item \emph{Lagerkosten} fallen proportional zum Lagerbestand an: Lagermiete, Energie, Versicherung, Kapital­bindung, Verderb. Sie steigen mit der Bestellmenge (weil größerer Durchschnittsbestand).
\end{itemize}

\subsection{Die Formel}

\begin{merksatz}[title={Andler-Formel}]
\[
x_{\text{opt}} = \sqrt{\frac{2 \cdot \text{Jahresbedarf} \cdot \text{Bestellkosten}}{\text{Einkaufspreis} \cdot \text{Lagerkostensatz}}}
\]
\end{merksatz}

\begin{eselsbruecke}
Im Zähler: \textbf{,,2 mal Bedarf mal Kosten''} (was raufgeht). Im Nenner:
\textbf{,,Preis mal Lagersatz''} (was Last erzeugt). Wurzel = Optimum.
\end{eselsbruecke}

\subsection{Beispielrechnung}

\begin{beispiel}
Ein Bürodienstleister benötigt pro Jahr 12\,000~Aktenordner. Die Bestellkosten pro
Bestellung betragen 40~€, der Einkaufspreis pro Ordner 2,50~€ und der Lagerkostensatz
20\,\% p.\,a.

\medskip
\[
x_{\text{opt}} = \sqrt{\frac{2 \cdot 12\,000 \cdot 40}{2{,}50 \cdot 0{,}20}}
= \sqrt{\frac{960\,000}{0{,}50}}
= \sqrt{1\,920\,000}
\approx \textbf{1\,386 Stück}
\]

\medskip
Die optimale Bestellmenge beträgt also rund 1\,386~Aktenordner pro Bestellung.
Bestellhäufigkeit: $12\,000 \div 1\,386 \approx 8{,}7$ Bestellungen pro Jahr ---
also etwa alle sechs Wochen.
\end{beispiel}

\subsection{Grenzen der Andler-Formel}

Die Andler-Formel ist eine Idealisierung. Sie nimmt an:

\begin{itemize}[itemsep=0pt]
  \item Konstanter Verbrauch über das Jahr (in Wirklichkeit: Saisonschwankungen)
  \item Konstanter Einkaufspreis (in Wirklichkeit: Mengenrabatte!)
  \item Konstanter Lagerkostensatz
  \item Sofortige Lieferung
  \item Keine Mindestbestellmengen
\end{itemize}

In der Praxis liefert sie deshalb einen \emph{Anhaltspunkt}, der dann an die realen
Bedingungen angepasst werden muss.

\section{Bestellung}

Eine schriftliche Bestellung enthält:

\begin{itemize}[itemsep=0pt]
  \item Eigene Adresse + Lieferadresse
  \item Datum + Bestellnummer (intern, zur Nachverfolgung)
  \item Genaue Artikelbezeichnung mit Artikelnummer
  \item Menge
  \item Preis (laut Angebot)
  \item Liefertermin
  \item Liefer- und Zahlungsbedingungen
  \item Bezug auf das vorhergehende Angebot (Datum/Nummer)
\end{itemize}

Form: schriftlich nach DIN~5008. Auch wenn ein mündlicher Vertrag rechtlich gültig
wäre --- die Schriftform schafft Beweissicherheit, falls es später zu Streit kommt.

\section{Bestellüberwachung}

\subsection{Was wird überwacht?}

Mit dem Versenden der Bestellung ist Ihre Arbeit nicht beendet. Bis die Ware im Lager
ist, müssen drei Aspekte überwacht werden:

\begin{itemize}[itemsep=0pt]
  \item \textbf{Termintreue}: Kommt die Lieferung pünktlich?
  \item \textbf{Mengentreue}: Stimmt die gelieferte Menge?
  \item \textbf{Qualität}: Ist die Ware mängelfrei? (das prüft die Wareneingangs­kontrolle, Tag~5)
\end{itemize}

\subsection{Mahnverfahren bei Lieferverzug}

Bleibt der Lieferant in Verzug, leitet man das kaufmännische Mahnverfahren ein.
Es kennt drei Stufen:

\begin{center}
\begin{tabular}{p{2cm} p{4cm} p{8cm}}
\toprule
\textbf{Stufe} & \textbf{Bezeichnung} & \textbf{Inhalt}\\
\midrule
1 & Erinnerung & Freundlicher Hinweis auf den Verzug, Bitte um Lieferung, neue Frist setzen\\
2 & Mahnung & Klarer Hinweis auf Verzug, Androhung von Schadensersatz\\
3 & Letzte Mahnung & Letzte Chance, Androhung von Rücktritt vom Vertrag, Schadensersatz und gerichtlichem Mahnverfahren\\
\bottomrule
\end{tabular}
\end{center}

Reagiert der Lieferant nach der dritten Mahnung nicht, wird entweder ein Inkassobüro
beauftragt oder das gerichtliche Mahnverfahren (Mahnbescheid beim zentralen Mahngericht)
eingeleitet.

\section{Nichtige und anfechtbare Rechtsgeschäfte (Vertiefung Tag 3)}

Im Zusammenhang mit Bestellungen lohnt ein kurzer Blick auf zwei Kategorien fehlerhafter
Verträge:

\subsection{Nichtige Rechtsgeschäfte}

Nichtige Rechtsgeschäfte sind \emph{von Anfang an} unwirksam --- sie entfalten nie
Rechts­wirkung. Gründe:

\begin{itemize}[itemsep=0pt]
  \item Geschäftsunfähigkeit (§\,105 BGB)
  \item Verstoß gegen ein Verbotsgesetz (§\,134)
  \item Sittenwidrigkeit oder Wucher (§\,138)
  \item Formmangel bei vorgeschriebener Form (z.\,B.\ Grundstück ohne Notar)
  \item Scheingeschäft (§\,117)
\end{itemize}

\subsection{Anfechtbare Rechtsgeschäfte}

Anfechtbare Rechtsgeschäfte sind \emph{zunächst wirksam}. Erst durch eine ausdrückliche
Anfechtungs­erklärung (§\,143 BGB) werden sie rückwirkend nichtig (§\,142). Anfechtungs­gründe:

\begin{itemize}[itemsep=0pt]
  \item \emph{Inhalts­irrtum} (§\,119 Abs.\,1) --- Sie wussten, was Sie sagten, aber irrten sich über die Bedeutung.
  \item \emph{Erklärungs­irrtum} (§\,119 Abs.\,1) --- Versprecher, Verschreiber.
  \item \emph{Eigenschafts­irrtum} (§\,119 Abs.\,2) --- Irrtum über eine wesentliche Eigenschaft der Sache oder Person.
  \item \emph{Arglistige Täuschung} (§\,123 Abs.\,1) --- Bewusste Irreführung durch den Vertragspartner.
  \item \emph{Drohung} (§\,123 Abs.\,1) --- Vertrag kam unter Druck zustande.
\end{itemize}

\begin{merksatz}
\textbf{Fristen}: Anfechtung wegen Irrtum unverzüglich nach Entdeckung (§\,121).
Anfechtung wegen Täuschung oder Drohung: 1~Jahr ab Entdeckung (§\,124). Nach 10~Jahren
in jedem Fall ausgeschlossen (§\,121 Abs.\,2, §\,124 Abs.\,3).
\end{merksatz}

\section{Zusammenfassung}

\begin{enumerate}
  \item Vier \textbf{Bestellverfahren}: Bestellpunkt, Bestellrhythmus, JIT, Kanban --- je nach Material und Lieferant.
  \item \textbf{Meldebestand} = Mindestbestand $+$ (Tagesverbrauch $\times$ Lieferzeit).
  \item \textbf{Andler-Formel} findet das Optimum zwischen Lager- und Bestellkosten.
  \item \textbf{Bestellung} schriftlich nach DIN 5008 (Beweissicherheit).
  \item \textbf{Bestellüberwachung}: Termin, Menge, Qualität.
  \item \textbf{Mahnverfahren}: 3 Stufen (Erinnerung, Mahnung, letzte Mahnung).
  \item \textbf{Anfechtung} bei Irrtum, Täuschung, Drohung --- macht den Vertrag rückwirkend nichtig.
\end{enumerate}

\newpage

\section*{Glossar}
\addcontentsline{toc}{section}{Glossar}

\glossareintrag{Bestellpunktverfahren}{Verbrauchsgesteuertes Bestellverfahren --- Bestellung wird ausgelöst, sobald der Lagerbestand den Meldebestand erreicht.}

\glossareintrag{Bestellrhythmusverfahren}{Zeitgesteuertes Verfahren --- Bestellung in festen Intervallen, Bestellmenge variabel.}

\glossareintrag{Mindestbestand}{Sicherheitspuffer im Lager; sollte nie unterschritten werden.}

\glossareintrag{Meldebestand}{Bestellauslösepunkt --- Mindestbestand plus erwarteter Verbrauch während der Lieferzeit.}

\glossareintrag{Höchstbestand}{Maximaler Lagerbestand nach einer neuen Lieferung; Mindestbestand $+$ Bestellmenge.}

\glossareintrag{Just-in-Time (JIT)}{Fertigungssynchrone Beschaffung; Lieferung exakt zum Bedarfszeit\-punkt, Lager praktisch null.}

\glossareintrag{Kanban}{Pull-Verfahren mit Karten oder Signalen; Material wird erst auf Anforderung des Verbrauchers geliefert. Aus Japan (Toyota).}

\glossareintrag{Andler-Formel}{Formel zur Berechnung der optimalen Bestellmenge: $x_\text{opt} = \sqrt{2 \cdot D \cdot K_b / (p \cdot s)}$.}

\glossareintrag{Bestellkosten}{Fixkosten pro Bestellvorgang (Personal, IT, Frachtsockel) --- sinken pro Stück bei größerer Bestellmenge.}

\glossareintrag{Lagerkosten}{Kosten proportional zum Lagerbestand (Miete, Energie, Versicherung, Kapital\-bindung) --- steigen mit größerer Bestellmenge.}

\glossareintrag{Lagerkostensatz}{Prozentualer Anteil der Lagerkosten am Lagerwert pro Jahr (typisch 10--25\,\%).}

\glossareintrag{Bestellüberwachung}{Kontrolle, ob die Lieferung termin- und mengentreu erfolgt.}

\glossareintrag{Mahnung}{Aufforderung an einen säumigen Schuldner, seine Pflicht zu erfüllen.}

\glossareintrag{Mahnstufen}{Drei Stufen des kaufmännischen Mahnverfahrens: Erinnerung, Mahnung, letzte Mahnung.}

\glossareintrag{Mahnbescheid}{Gerichtliches Mahnverfahren beim zentralen Mahngericht --- führt nach Schweigen zum Vollstreckungsbescheid.}

\glossareintrag{Nichtiges Rechtsgeschäft}{Vertrag, der von Anfang an unwirksam ist (z.\,B.\ wegen Geschäfts\-unfähigkeit, Sittenwidrigkeit).}

\glossareintrag{Anfechtbares Rechtsgeschäft}{Vertrag, der zunächst wirksam ist und durch Anfechtungserklärung rückwirkend nichtig wird.}

\glossareintrag{Anfechtungsgründe}{Inhalts\-irrtum (§\,119), Eigenschafts\-irrtum, arglistige Täuschung (§\,123), Drohung (§\,123).}

\glossareintrag{Anfechtungsfrist}{Bei Irrtum unverzüglich nach Entdeckung; bei Täuschung/Drohung 1~Jahr; in jedem Fall max.\ 10~Jahre.}

\end{document}
